\subsection{Algebraic elements of an algebra}

Let \(\mathbf{A}\) be an algebra over a field \(\mathbf{K}\) and \(x\) an element of \(\mathbf{A}\) . Two cases are possible:

\begin{enumerate}
\def\labelenumi{\alph{enumi})}
\item
  The family of monomials \((x^{n})_{n\in\mathbb{N}}\) is free over K. Then we say that x is transcendental over K. There is an isomorphism of the polynomial algebra K{[}X{]} onto the subalgebra K{[}x{]} of A generated by x, and the latter is of infinite degree over K.
\item
  There exists an integer \(n \geqslant 1\) such that the monomials \(1, x, \ldots, x^{n-1}, x^n\) are linearly dependent; it comes to the same to say that there exists a polynomial \(f \neq 0\) in \(\mathbf{K}[\mathbf{X}]\) such that \(f(x) = 0\) . Then we say that \(x\) is algebraic over \(\mathbf{K}\) . The least integer \(n \geqslant 1\) satisfying the above property is called the degree of \(x\) over \(\mathbf{K}\) . If the degree of \(x\) over \(\mathbf{K}\) is \(n\) , then the monomials \(1, x, \ldots, x^{n-1}\) are linearly independent over \(\mathbf{K}\) and there exist elements \(a_0, a_1, \ldots, a_{n-1}\) of \(\mathbf{K}\) such that
\end{enumerate}

\[
x ^ {n} = a _ {0} + a _ {1} x + \dots + a _ {n - 1} x ^ {n - 1}.
\]

The polynomial \(f(\mathbf{X}) = \mathbf{X}^n - \sum_{k=0}^{n-1} a_k \mathbf{X}^k\) is the unique monic polynomial of degree

\(n\) in \(\mathbf{K}[\mathbf{X}]\) such that \(f(x) = 0\) ; it is called the minimal polynomial of \(x\) over \(\mathbf{K}\) .

THEOREM 1. --- Let A be an algebra over a field K, x an element of A algebraic over K, n the degree and f the minimal polynomial of x over K.

\begin{enumerate}
\def\labelenumi{\alph{enumi})}
\item
  For a polynomial \(g \in \mathbf{K}[\mathbf{X}]\) to satisfy \(g(x) = 0\) it is necessary and sufficient that \(g\) should be a multiple of \(f\) .
\item
  The mapping \(g \mapsto g(x)\) defines by passage to quotients an isomorphism of the quotient algebra \(\mathbf{K}[\mathbf{X}] / (f)\) onto the algebra \(\mathbf{K}[x]\) , and the elements 1, x, \ldots, \(x^{n-1}\) form a basis of \(\mathbf{K}[x]\) over \(\mathbf{K}\) . In particular, \([\mathbf{K}[x]: \mathbf{K}] = n\) .
\item
  Assume that A is an integral domain. Then K{[}x{]} is a field and f the unique monic irreducible polynomial in K{[}X{]} such that \(f(x) = 0\) .
\item
  For \(x\) to be invertible in \(\mathbf{A}\) it is necessary and sufficient that \(f(0) \neq 0\) ; then we have \(x^{-1} \in \mathbf{K}[x]\) .
\end{enumerate}

There exists a unique algebra homomorphism \(\varphi: \mathbf{K}[\mathbf{X}] \to \mathbf{A}\) such that \(\varphi(\mathbf{X}) = x\) ; we have \(\varphi(\mathbf{P}) = \mathbf{P}(x)\) for every \(\mathbf{P} \in \mathbf{K}[\mathbf{X}]\) and the image of \(\varphi\) equals \(\mathbf{K}[x]\) . Let \(\mathfrak{a}\) be the kernel of \(\varphi\) ; by construction, the minimal polynomial \(f\) of \(x\) over \(\mathbf{K}\) belongs to \(\mathfrak{a}\) and it is the monic polynomial of least degree in \(\mathfrak{a}\) . Therefore (IV, p.~11, Prop. 11) we have \(\mathfrak{a} = (f)\) , whence \(a\) ). Assertion \(b\) ) follows at once from \(a\) ) and the Cor. of Prop. 10 of IV, p.~11.

Suppose that A is an integral domain. The algebra K{[}x{]} is then an integral domain and of finite degree over K, hence it is a field (V, p.~10, Cor.). The ideal (f) of K{[}x{]} is thus maximal, that means that f is irreducible in K{[}X{]} (IV, p.~13). Finally, let g be a monic irreducible polynomial in K{[}X{]} such that \(g(x) = 0\) ; by a) it is a multiple of f, hence g = f, and this proves c).

It remains to prove \(d\) ). There exists a polynomial \(g \in \mathbf{K}[\mathbf{X}]\) of degree \(n - 1\) and an element \(a\) of \(\mathbf{K}\) such that \(f(\mathbf{X}) = \mathbf{X}g(\mathbf{X}) + a\) , whence \(f(0) = a\) . If \(a = 0\) , we have \(xg(x) = f(x) = 0\) , and \(g(x) \neq 0\) , so then \(x\) is not invertible in \(\mathbf{A}\) . If on the other hand \(a \neq 0\) , then we have \(x \cdot [-a^{-1}g(x)] = 1\) , so \(x\) is then invertible in \(\mathbf{A}\) and \(x^{-1} = -a^{-1}g(x)\) .

COROLLARY 1. --- Let A be an algebra over a field K. For an element x of A to be algebraic over K it is necessary and sufficient that the subalgebra \(K[x]\) of A generated by x should be of finite degree over K. In particular, if A is of finite degree over K, every element of A is algebraic over K.

COROLLARY 2. --- Let E be an extension of K, A an algebra over E and x an element of A algebraic over K. Then x is algebraic over E, the minimal polynomial of x over E divides the minimal polynomial of x over K and the degree of x over E is at most equal to the degree of x over K.

For let \(f\) be the minimal polynomial of \(x\) over \(\mathbf{K}\) ; we have \(f(x) = 0\) and \(f \in \operatorname{E}[\mathbf{X}]\) , hence \(x\) is algebraic over \(\mathbf{E}\) and \(f\) is a multiple of the minimal polynomial of \(x\) over \(\mathbf{E}\) (Th. 1, a)).

Remark. --- Let E be an extension of a field K and x an element of E which is root of a monic irreducible polynomial \(f \in K[X]\) . Th. 1, c) then shows that f is the minimal polynomial of x over K.

Examples. --- * 1) In the field of complex numbers C, the number i is algebraic of degree 2 over the prime field Q; for if \(f(\mathbf{X}) = \mathbf{X}^{2} + 1\) , then \(f(i) = 0\) , and \(x^{2} + 1 \neq 0\) for all \(x \in Q\) , hence \(i \notin Q\) . The field \(\mathbf{Q}(i)\) is thus an extension of degree 2 of Q; it consists of all numbers \(a + bi\) , where a, b are rational. Likewise i is algebraic of degree 2 over the field R of real numbers, and C is an extension of degree 2 of R.

\begin{enumerate}
\def\labelenumi{\arabic{enumi})}
\setcounter{enumi}{1}
\tightlist
\item
  Let K be a field and F the field \(\mathrm{K}(\mathrm{X})\) of rational functions in one indeterminate over K. Let E be the subfield \(\mathrm{K}(\mathrm{X}^{3})\) of F; we have \(\mathrm{F} = \mathrm{E}(\mathrm{X})\) and X is algebraic over E, since it is a root of the polynomial \(Y^{3} - X^{3}\) of the ring \(E[Y]\) ; this polynomial is irreducible in \(E[Y]\) , for in the contrary case it would have at least one factor of the first degree, and there would then exist two non-zero polynomials \(u(\mathrm{X})\) , \(v(\mathrm{X})\) of \(K[X]\) such that \((u(\mathrm{X}^{3}))^{3} = \mathrm{X}^{3}(v(\mathrm{X}^{3}))^{3}\) which is absurd, for if m and n are the degrees of u and v, this would imply \(9m = 9n + 3\) , or \(3m = 3n + 1\) . The field F is thus of degree 3 over E, and every element of F may be written in just one way as a linear combination \(f(\mathrm{X}^{3}) + \mathrm{X}g(\mathrm{X}^{3}) + \mathrm{X}^{2}h(\mathrm{X}^{3})\) , where f, g, h are three rational fractions of \(\mathrm{K}(\mathrm{X})\) .
\end{enumerate}

*3) In the field \(\mathbf{R}\) of real numbers it may be shown \(^1\) that the number \(\pi\) is transcendental over the prime field \(\mathbf{Q}\) . *

